\subsection{Divisibility relations in a field}

Here we will define certain ordered groups which play an important part in algebra. The notation usually used for these groups is multiplicative; the application of results previously obtained in additive notation to these groups thus presupposes the translation to multiplicative notation --- which should present no difficulty to the reader. Throughout this section, the letter A will denote an integral domain and K its field of fractions (I, p.~116).

In the multiplicative group \(\mathbf{K}^*\) of nonzero elements of \(\mathbf{K}\) , the set \(\mathbf{P}\) of nonzero elements of \(\mathbf{A}\) is closed under multiplication, since \(\mathbf{A}\) is a ring. Hence it defines a pre-order relation on \(\mathbf{K}^*\) by \(x^{-1}y\in \mathbf{P}\) , that is « there exists \(z\in \mathbf{P}\) such that \(y = zx\) », which makes \(\mathbf{K}^*\) a pre-ordered group (written multiplicatively) (VI, p.~3, Prop. 3). Extending the terminology relating to elements of \(\mathbf{A}\) (I, p.~97) to \(\mathbf{K}^*\) , the relation \(x^{-1}y\in \mathbf{P}\) can also be expressed as: \(x\) divides \(y\) , or \(x\) is a divisor of \(y\) , or \(y\) is a multiple of \(x\) (relative to the ring \(\mathbf{A}\) ); and we will call the relation \(x^{-1}y\in \mathbf{P}\) the divisibility relation in \(\mathbf{K}^*\) relative to the ring \(\mathbf{A}\) . The relation « \(x\) divides \(y\) » is denoted \(x|y\) , and its negation \(x\nmid y\) . The elements of \(\mathbf{P}\) are none other than the multiples of 1.

Remarks. --- 1) The divisibility relation in \(\mathbf{K}^*\) depends in an essential way on the particular ring \(\mathbf{A}\) . If \(\mathbf{A} = \mathbf{K}\) we obtain the « trivial » relation, under which \(x \mid y\) for every pair \((x, y)\) of elements of \(\mathbf{K}^*\) . Let \(p\) and \(q\) be prime numbers; the rational numbers \(r/s\) whose denominators are not multiples of \(p\) (resp. \(q\) ) form a subring \(\mathbf{Z}_{(p)}\) (resp. \(\mathbf{Z}_{(q)}\)) of \(\mathbf{Q}\) ; the divisibility relations in \(\mathbf{Q}^*\) relative to these two rings are distinct if \(p \neq q\) , the number \(p/q\) being a multiple of 1 under one relation but not under the other.

\begin{enumerate}
\def\labelenumi{\arabic{enumi})}
\setcounter{enumi}{1}
\tightlist
\item
  We will sometimes extend the definition of the relation \(x \mid y\) to pairs of elements of \(K\) (rather than only \(K^*\) ), this relation being taken to mean « there exists \(z \in A\) such that \(y = zx\) »; hence we will have \(x \mid 0\) for all \(x \in K\) . This allows us to state the following results without restrictions: if \(x \mid y\) and \(x \mid z\) then \(x \mid (y - z)\) ; if \(x \mid y\) and \(x \nmid z\) then \(x \nmid (y - z)\) . In the same way we can extend all the corresponding terminology.
\end{enumerate}

To obtain an order relation (No.~2) from the divisibility relation we have to pass to the quotient group of \(\mathbf{K}^*\) by the subgroup \(\mathbf{A}^*\) of elements \(x \in \mathbf{K}^*\) such that \(x \mid 1\) and \(1 \mid x\) ; these are the elements of \(\mathbf{P}\) which are divisors of 1, that is to say the invertible elements of \(\mathbf{A}\) ; by abuse of language they are often called the units of the ring \(\mathbf{A}\) . The quotient group \(\mathbf{K}^* / \mathbf{A}^*\) is then an ordered group. Two elements \(x\) and \(y\) of \(\mathbf{K}^*\) which belong to the same coset of \(\mathbf{A}^*\) are said to be associate; this means that \(x \mid y\) and \(y \mid x\) . If on the other hand \(x\) divides \(y\) without \(y\) dividing \(x\) , we say that \(x\) strictly divides \(y\) , or that \(x\) is a strict divisor of \(y\) , or that \(y\) is a strict multiple of \(x\) .

Note that \(K^{*}/A^{*}\) is a filtered group, since K is the field of fractions of A (VI, p.~4, Prop. 4).

To say that two elements x and y of \(K^{*}\) are associate amounts to saying that they have the same multiples in K, by virtue of the transitivity of the divisibility relation. For all \(x \in K\) , we will denote by Ax the set of all elements zx with \(z \in A\) ; the set Ax is a submodule of K regarded as an A-module. Extending the terminology relating to the case where \(x \in A\) , we will call it a principal fractional ideal of the field K relative to the ring A. In contrast the ideals of the ring A are said to be integral.

\pandocbounded{\includegraphics[keepaspectratio]{images/part002_d9406660.jpg}}

Note that if \(\mathbf{A} \neq \mathbf{K}\) then a principal fractional ideal \(\neq \{0\}\) is not an ideal of \(\mathbf{K}\) considered as a ring.

The principal fractional ideal Ax is also denoted \((x)\) . We will write \(x \equiv 0 \pmod{y}\) for \(x \in Ay\) and \(x \equiv x' \pmod{y}\) for \(x - x' \in Ay\) ; if \(x \equiv x' \pmod{y}\) then \(zx \equiv zx' \pmod{zy}\) for any \(z \in K\) .

Note that \(x \equiv x' (\bmod y)\) does not imply \(zx \equiv zx' (\bmod y)\) unless \(z \in A\) . Thus in \(\mathbf{Q}\) , relative to \(\mathbf{Z}\) , we have \(4 \equiv 2 (\bmod 2)\) but not \(2 \equiv 1 (\bmod 2)\) .

The relation \(x \mid y\) is obviously equivalent to \((x) \supset (y)\) . The mapping \(x \mapsto (x)\) from \(K^*\) onto the set \(\mathcal{P}^*\) of principal fractional ideals \(\neq (0)\) of \(K\) thus defines, on passing to the quotient, a bijective mapping from \(K^*/A^*\) onto \(\mathcal{P}^*\) ; translating the group structure of \(K^*/A^*\) to \(\mathcal{P}^*\) by means of this mapping, we are led to define the product of two principal fractional ideals \((x)\) and \((y)\) to be the ideal \((xy)\) , which depends only on \((x)\) and \((y)\) . Equipped with this law and the order relation \((x) \supset (y)\) , the set \(\mathcal{P}^*\) is an ordered group, isomorphic to \(K^*/A^*\) . By convention we will identify \(\mathcal{P}^*\) with \(K^*/A^*\) via the above map.

Note that the relation « x divides y » which, in the case of positive integers, implies that x is smaller than y, corresponds to the inclusion \((x) \supset (y)\) , in which the ideal \((x)\) is « greater » than the ideal \((y)\) . We can keep this « order reversal » in mind by noting that for example 7 has « more multiples » than 91.

If we extend the relation \(x \mid y\) to all elements of \(\mathbf{K}\) , this relation is still equivalent to \((x) \supset (y)\) in the set \(\mathcal{P}\) of all principal fractional ideals of \(\mathbf{K}\) (in which (0) is the smallest element under the relation of inclusion).

As in the previous sections, we will generally be using additive notation in the sequel. However, terminology relating to divisibility will be introduced following the corresponding additive terminology, in paragraphs preceded by the sign (DIV) (in which it is understood that the notation used is that of the present section). In order to make the reader's task easier, certain results will be translated into the language of divisibility, the translation of Prop. 7, for example, being denoted « PROPOSITION 7 (DIV) ».

